% Worked Examples: Topic 2.6 - Frequency Domain Analysis
\documentclass[11pt,a4paper]{article}
\usepackage{../worked-examples-template}

\title{\textbf{Worked Examples}\\[0.5em]
\large Topic 2.6: Frequency Domain Analysis\\[0.3em]
\normalsize \coursesubtitle{} -- \coursetitle}
\author{Assoc.\ Prof.\ William Robertson \and Dr Sean McGowan}
\date{\today}

\begin{document}

\maketitle
\tableofcontents
\newpage

\section{Concept 2.6.1: Loop Analysis}

\subsection{Example 1: Closed-Loop System Analysis}

\paragraph{Problem:} Consider a unity feedback system with open-loop transfer function:
\begin{align}
L(s) = C(s)P(s) = \frac{K}{s(s+2)}
\end{align}

Find: (a) the closed-loop transfer function, (b) the value of $K$ that gives closed-loop poles with damping ratio $\zeta = 0.707$, and (c) the corresponding natural frequency $\omega_0$.

\paragraph{Solution:}

\textbf{Part (a): Closed-Loop Transfer Function}

For unity feedback, the closed-loop transfer function is:
\begin{align}
T(s) = \frac{L(s)}{1 + L(s)} = \frac{K/[s(s+2)]}{1 + K/[s(s+2)]} = \frac{K}{s^2 + 2s + K}
\end{align}

\textbf{Part (b) \& (c): Finding $K$ for $\zeta = 0.707$}

The characteristic equation is:
\begin{align}
s^2 + 2s + K = 0
\end{align}

Compare with the standard second-order form:
\begin{align}
s^2 + 2\zeta\omega_0 s + \omega_0^2 = 0
\end{align}

Matching coefficients:
\begin{align}
2\zeta\omega_0 &= 2 \\
\omega_0^2 &= K
\end{align}

From the first equation with $\zeta = 0.707$:
\begin{align}
\omega_0 = \frac{2}{2\zeta} = \frac{2}{2(0.707)} = \frac{1}{0.707} \approx 1.414 \text{ rad/s}
\end{align}

Therefore:
\begin{align}
K = \omega_0^2 = (1.414)^2 \approx 2
\end{align}

\textbf{Result:}
\begin{itemize}
\item Closed-loop transfer function: $T(s) = \frac{K}{s^2 + 2s + K}$
\item For $\zeta = 0.707$: $K = 2$ and $\omega_0 = \sqrt{2} \approx 1.414$ rad/s
\item The closed-loop poles are at: $s = -\zeta\omega_0 \pm j\omega_0\sqrt{1-\zeta^2} = -1 \pm j$
\end{itemize}

\textbf{Verification:} With $K = 2$, the characteristic equation is $s^2 + 2s + 2 = 0$:
\begin{align}
s = \frac{-2 \pm \sqrt{4-8}}{2} = \frac{-2 \pm j2}{2} = -1 \pm j
\end{align}

Damping ratio: $\zeta = \cos(\arctan(1/1)) = \cos(45°) = 0.707$ ✓

\section{Concept 2.6.2: Nyquist Criterion}

\subsection{Example 1: Applying the Nyquist Stability Criterion}

\paragraph{Problem:} Use the Nyquist criterion to determine the stability of a unity feedback system with open-loop transfer function:
\begin{align}
L(s) = \frac{K}{(s+1)(s+2)}
\end{align}

for $K = 5$.

\paragraph{Solution:}

The Nyquist stability criterion states that for a closed-loop system:
\begin{align}
N = P - Z
\end{align}

where:
\begin{itemize}
\item $N$ = number of clockwise encirclements of the point $-1 + j0$ by the Nyquist plot
\item $P$ = number of open-loop poles in the right half-plane
\item $Z$ = number of closed-loop poles in the right half-plane
\end{itemize}

For stability, we need $Z = 0$, which requires $N = P$.

\textbf{Step 1: Count open-loop RHP poles}

The open-loop poles are at $s = -1$ and $s = -2$ (both in left half-plane).
Therefore: $P = 0$

\textbf{Step 2: Analyze the Nyquist plot}

For the frequency response $L(j\omega)$ with $K = 5$:
\begin{align}
L(j\omega) = \frac{5}{(j\omega+1)(j\omega+2)} = \frac{5}{(1+j\omega)(2+j\omega)}
\end{align}

Key points:
\begin{itemize}
\item At $\omega = 0$: $L(j0) = \frac{5}{1 \cdot 2} = 2.5$ (on positive real axis)
\item As $\omega \to \infty$: $L(j\omega) \to 0$ (magnitude approaches 0)
\item Phase is always negative (two poles, no zeros)
\end{itemize}

To check if the plot crosses the negative real axis at $-1$:

At crossover, the imaginary part is zero and real part is negative. For $L(j\omega)$ to be real and negative:
\begin{align}
\text{Im}[L(j\omega)] = 0
\end{align}

Converting to magnitude and phase:
\begin{align}
|L(j\omega)| &= \frac{5}{\sqrt{(1+\omega^2)(4+\omega^2)}} \\
\angle L(j\omega) &= -\arctan(\omega) - \arctan(\omega/2)
\end{align}

The phase crosses $-180°$ when:
\begin{align}
\arctan(\omega) + \arctan(\omega/2) = 180° = \pi
\end{align}

This never happens for finite $\omega$ (only as $\omega \to \infty$, where magnitude is 0).

The Nyquist plot:
\begin{itemize}
\item Starts at 2.5 on the positive real axis ($\omega = 0$)
\item Curves through the lower half-plane (negative imaginary part)
\item Approaches origin as $\omega \to \infty$
\item Never encircles $-1 + j0$
\end{itemize}

Therefore: $N = 0$

\textbf{Step 3: Apply Nyquist criterion}

\begin{align}
Z = P - N = 0 - 0 = 0
\end{align}

\textbf{Result:} The closed-loop system is stable for $K = 5$ since there are no closed-loop poles in the right half-plane ($Z = 0$).

\section{Concept 2.6.3: Stability Margins}

\subsection{Example 1: Calculating Gain and Phase Margins}

\paragraph{Problem:} For the open-loop transfer function:
\begin{align}
L(s) = \frac{10}{s(s+1)(s+5)}
\end{align}

Calculate: (a) the gain crossover frequency $\omega_{gc}$ and phase margin, and (b) the phase crossover frequency $\omega_{pc}$ and gain margin.

\paragraph{Solution:}

\textbf{Part (a): Gain Crossover Frequency and Phase Margin}

The gain crossover frequency $\omega_{gc}$ is where $|L(j\omega_{gc})| = 1$ (or 0 dB).

\begin{align}
|L(j\omega)| = \frac{10}{|\omega|\sqrt{1+\omega^2}\sqrt{25+\omega^2}}
\end{align}

Setting $|L(j\omega_{gc})| = 1$:
\begin{align}
\frac{10}{\omega_{gc}\sqrt{1+\omega_{gc}^2}\sqrt{25+\omega_{gc}^2}} = 1
\end{align}

This is transcendental and typically solved numerically. For approximation, try $\omega_{gc} \approx 1.4$ rad/s:
\begin{align}
|L(j1.4)| = \frac{10}{1.4\sqrt{1+1.96}\sqrt{25+1.96}} = \frac{10}{1.4 \times 1.72 \times 5.19} \approx 0.8
\end{align}

Try $\omega_{gc} \approx 1.2$ rad/s:
\begin{align}
|L(j1.2)| = \frac{10}{1.2\sqrt{1+1.44}\sqrt{25+1.44}} = \frac{10}{1.2 \times 1.56 \times 5.14} \approx 1.04 \approx 1
\end{align}

So $\omega_{gc} \approx 1.2$ rad/s.

The phase margin is:
\begin{align}
PM = 180° + \angle L(j\omega_{gc})
\end{align}

where:
\begin{align}
\angle L(j\omega) = -90° - \arctan(\omega) - \arctan(\omega/5)
\end{align}

At $\omega = 1.2$:
\begin{align}
\angle L(j1.2) &= -90° - \arctan(1.2) - \arctan(0.24) \\
&= -90° - 50.2° - 13.5° \\
&= -153.7°
\end{align}

Therefore:
\begin{align}
PM = 180° - 153.7° = 26.3°
\end{align}

\textbf{Part (b): Phase Crossover Frequency and Gain Margin}

The phase crossover frequency $\omega_{pc}$ is where $\angle L(j\omega_{pc}) = -180°$.

\begin{align}
-90° - \arctan(\omega_{pc}) - \arctan(\omega_{pc}/5) = -180°
\end{align}

Simplifying:
\begin{align}
\arctan(\omega_{pc}) + \arctan(\omega_{pc}/5) = 90°
\end{align}

Using the tangent addition formula: $\tan(A + B) = \frac{\tan A + \tan B}{1 - \tan A \tan B}$

Let $A = \arctan(\omega_{pc})$ and $B = \arctan(\omega_{pc}/5)$:
\begin{align}
\tan(90°) = \infty = \frac{\omega_{pc} + \omega_{pc}/5}{1 - \omega_{pc} \cdot \omega_{pc}/5}
\end{align}

This requires the denominator to be zero:
\begin{align}
1 - \frac{\omega_{pc}^2}{5} = 0 \quad \Rightarrow \quad \omega_{pc} = \sqrt{5} \approx 2.24 \text{ rad/s}
\end{align}

The gain margin is:
\begin{align}
GM = \frac{1}{|L(j\omega_{pc})|} \text{ (linear)} \quad \text{or} \quad GM_{dB} = -20\log_{10}|L(j\omega_{pc})|
\end{align}

At $\omega_{pc} = \sqrt{5}$:
\begin{align}
|L(j\sqrt{5})| = \frac{10}{\sqrt{5}\sqrt{1+5}\sqrt{25+5}} = \frac{10}{\sqrt{5}\sqrt{6}\sqrt{30}} = \frac{10}{\sqrt{900}} = \frac{10}{30} = \frac{1}{3}
\end{align}

Therefore:
\begin{align}
GM = \frac{1}{1/3} = 3 \text{ (linear)} = 20\log_{10}(3) \approx 9.5 \text{ dB}
\end{align}

\textbf{Results:}
\begin{itemize}
\item Gain crossover frequency: $\omega_{gc} \approx 1.2$ rad/s
\item Phase margin: $PM \approx 26°$ (marginally stable)
\item Phase crossover frequency: $\omega_{pc} = \sqrt{5} \approx 2.24$ rad/s
\item Gain margin: $GM = 3$ (or 9.5 dB)
\end{itemize}

\section{Concept 2.6.4: Minimum Phase Systems}

\subsection{Example 1: Identifying Minimum Phase Systems}

\paragraph{Problem:} Determine which of the following transfer functions are minimum phase systems:

(a) $G_1(s) = \frac{s+2}{(s+1)(s+3)}$

(b) $G_2(s) = \frac{s-2}{(s+1)(s+3)}$

(c) $G_3(s) = \frac{2e^{-0.5s}}{s+1}$

For each system, explain your reasoning.

\paragraph{Solution:}

A system is \textbf{minimum phase} if all its poles and zeros are in the left half of the complex plane (i.e., have negative real parts), and it has no time delays.

\textbf{Part (a): $G_1(s) = \frac{s+2}{(s+1)(s+3)}$}

Zeros: $s = -2$ (left half-plane) ✓

Poles: $s = -1, -3$ (both left half-plane) ✓

No time delay ✓

\textbf{Result:} $G_1(s)$ is a \textbf{minimum phase system}.

\textbf{Part (b): $G_2(s) = \frac{s-2}{(s+1)(s+3)}$}

Zeros: $s = +2$ (right half-plane) ✗

Poles: $s = -1, -3$ (both left half-plane) ✓

No time delay ✓

\textbf{Result:} $G_2(s)$ is \textbf{NOT minimum phase} because it has a right half-plane zero (at $s = +2$). This is called a \textit{non-minimum phase zero}.

\textbf{Part (c): $G_3(s) = \frac{2e^{-0.5s}}{s+1}$}

Zeros: None (the time delay $e^{-0.5s}$ is not a zero in the traditional sense)

Poles: $s = -1$ (left half-plane) ✓

Has time delay: $e^{-0.5s}$ ✗

\textbf{Result:} $G_3(s)$ is \textbf{NOT minimum phase} because it contains a time delay. Time delays introduce infinite phase lag without affecting magnitude, making the system non-minimum phase.

\textbf{Summary Table:}

\begin{center}
\begin{tabular}{|c|c|c|c|}
\hline
System & LHP Zeros & LHP Poles & No Delay & Minimum Phase? \\
\hline
$G_1(s)$ & ✓ & ✓ & ✓ & Yes \\
$G_2(s)$ & ✗ (RHP zero) & ✓ & ✓ & No \\
$G_3(s)$ & ✓ & ✓ & ✗ (has delay) & No \\
\hline
\end{tabular}
\end{center}

\textbf{Key Properties of Minimum Phase Systems:}
\begin{itemize}
\item Unique relationship between magnitude and phase response
\item Phase can be reconstructed from magnitude (and vice versa)
\item Minimum possible phase lag for a given magnitude response
\item Generally easier to control than non-minimum phase systems
\end{itemize}

\textbf{Practical Implications:}
\begin{itemize}
\item Non-minimum phase zeros (RHP zeros) impose fundamental limitations on achievable bandwidth and performance
\item Time delays make control more challenging and limit achievable phase margin
\item Example: inverted pendulum has non-minimum phase dynamics
\end{itemize}

\end{document}
